% GENERATED FILE. Edit canonical lessons, then run npm run generate:books.
% canonical-source-hash: fnv1a64:8b3547037216b7a5
% canonical-lessons: KA-C382-sangati, KA-C382-vastava, KA-C382-bhrame, KA-C382-astitva, KA-C382-managanu

\chapter{Facts and Illusions}
\label{ch:ka-facts-and-illusions}
\hlchaptermodality{382}

\begin{chapteropening}
\textbf{By the end of this chapter:} \emph{I can say a fact, reality, an illusion, existence and to realise.}

Complete the last lesson of chapter 382: I can say a fact, reality, an illusion, existence and to realise.
\end{chapteropening}

\section[\texorpdfstring{\kn{ಸಂಗತಿ} (saṅgati)}{ಸಂಗತಿ (saṅgati)}]{\kn{ಸಂಗತಿ} (saṅgati) — a fact; a matter}

\label{lesson:KA-C382-sangati}



\begin{quote}
Before the new one: say the Kannada for to reflect on, then the Kannada for intellect.
\end{quote}

\subsection*{You'll want to know: \kn{ಸಂಗತಿ}}
\textbf{\kn{ಸಂಗತಿ}} — \emph{saṅgati} — \textquotedblleft{}a fact; a matter\textquotedblright{}.

\begin{quote}
\kn{ಈ} \kn{ಸಂಗತಿ} \kn{ಯಾರಿಗೂ} \kn{ಗೊತ್ತಿರಲಿಲ್ಲ}.
\end{quote}

\begin{quote}
\emph{ī saṅgati yārigū gottiralilla.}
\end{quote}

\begin{quote}
Nobody knew this fact.
\end{quote}

\kn{ವಿಷಯ} is a topic; \kn{ಸಂಗತಿ} is something that is so.

\begin{cousinweb}
\textbf{\kn{ಸಂಗತಿ}} (\emph{saṅgati}) is borrowed from Sanskrit \emph{saṃgati}, \textquotedblleft{}a coming together, a meeting\textquotedblright{} (Monier-Williams); Kannada uses it for a matter or a fact.
\end{cousinweb}

\begin{grammarlens}[title={What you've built}]
One more word for giving reasons and opinions.
\end{grammarlens}

\subsection*{Your turn}
\begin{itemize}
\raggedright
  \item \emph{Say it:} \emph{saṅgati}
  \item \emph{Say it:} \emph{saṅgati}, once more
  \item \emph{Say it:} \emph{ī saṅgati yārigū gottiralilla.} — Nobody knew this fact.
  \item \emph{Recall:} say the Kannada for to reflect on, then the Kannada for intellect, then say \emph{saṅgati} again
\end{itemize}

\begin{tcolorbox}[breakable,colback=teal!4,colframe=teal!35!black,title={Before you move on}]
What does \kn{ಸಂಗತಿ} mean? (\textquotedblleft{}A fact; a matter\textquotedblright{}.) Say it once more.
\end{tcolorbox}
\section[\texorpdfstring{\kn{ವಾಸ್ತವ} (vāstava)}{ವಾಸ್ತವ (vāstava)}]{\kn{ವಾಸ್ತವ} (vāstava) — reality}

\label{lesson:KA-C382-vastava}



\begin{quote}
Before the new one: say the Kannada for intellect, then the Kannada for a fact.
\end{quote}

\subsection*{You'll want to know: \kn{ವಾಸ್ತವ}}
\textbf{\kn{ವಾಸ್ತವ}} — \emph{vāstava} — \textquotedblleft{}reality\textquotedblright{}.

\begin{quote}
\kn{ಕನಸು} \kn{ಬೇರೆ}, \kn{ವಾಸ್ತವ} \kn{ಬೇರೆ}.
\end{quote}

\begin{quote}
\emph{kanasu bēre, vāstava bēre.}
\end{quote}

\begin{quote}
Dreams are one thing, reality another.
\end{quote}

\kn{ವಾಸ್ತವವಾಗಿ}, in fact.

\begin{cousinweb}
\textbf{\kn{ವಾಸ್ತವ}} (\emph{vāstava}) is borrowed from Sanskrit \emph{vāstava}, \textquotedblleft{}substantial, real, true\textquotedblright{} (Monier-Williams), from \emph{vastu}, \textquotedblleft{}a thing\textquotedblright{}.
\end{cousinweb}

\begin{grammarlens}[title={What you've built}]
One more word for giving reasons and opinions.
\end{grammarlens}

\subsection*{Your turn}
\begin{itemize}
\raggedright
  \item \emph{Say it:} \emph{vāstava}
  \item \emph{Say it:} \emph{vāstava}, once more
  \item \emph{Say it:} \emph{kanasu bēre, vāstava bēre.} — Dreams are one thing, reality another.
  \item \emph{Recall:} say the Kannada for intellect, then the Kannada for a fact, then say \emph{vāstava} again
\end{itemize}

\begin{tcolorbox}[breakable,colback=teal!4,colframe=teal!35!black,title={Before you move on}]
What does \kn{ವಾಸ್ತವ} mean? (\textquotedblleft{}Reality\textquotedblright{}.) Say it once more.
\end{tcolorbox}
\section[\texorpdfstring{\kn{ಭ್ರಮೆ} (bhrame)}{ಭ್ರಮೆ (bhrame)}]{\kn{ಭ್ರಮೆ} (bhrame) — an illusion}

\label{lesson:KA-C382-bhrame}



\begin{quote}
Before the new one: say the Kannada for a fact, then the Kannada for reality.
\end{quote}

\subsection*{You'll want to know: \kn{ಭ್ರಮೆ}}
\textbf{\kn{ಭ್ರಮೆ}} — \emph{bhrame} — \textquotedblleft{}an illusion\textquotedblright{}.

\begin{quote}
\kn{ನೀರು} \kn{ಇದೆ} \kn{ಅನ್ನುವುದು} \kn{ಬರೀ} \kn{ಭ್ರಮೆ} \kn{ಆಗಿತ್ತು}.
\end{quote}

\begin{quote}
\emph{nīru ide annuvudu barī bhrame āgittu.}
\end{quote}

\begin{quote}
That there was water was only an illusion.
\end{quote}

\kn{ಭ್ರಮೆಯಲ್ಲಿ} \kn{ಇರು}, to be under an illusion.

\begin{cousinweb}
\textbf{\kn{ಭ್ರಮೆ}} (\emph{bhrame}) is borrowed from Sanskrit \emph{bhrama}, \textquotedblleft{}wandering about\textquotedblright{}, and so \textquotedblleft{}confusion, error\textquotedblright{} (CDIAL 9646).
\end{cousinweb}

\begin{grammarlens}[title={What you've built}]
One more word for giving reasons and opinions.
\end{grammarlens}

\subsection*{Your turn}
\begin{itemize}
\raggedright
  \item \emph{Say it:} \emph{bhrame}
  \item \emph{Say it:} \emph{bhrame}, once more
  \item \emph{Say it:} \emph{nīru ide annuvudu barī bhrame āgittu.} — That there was water was only an illusion.
  \item \emph{Recall:} say the Kannada for a fact, then the Kannada for reality, then say \emph{bhrame} again
\end{itemize}

\begin{tcolorbox}[breakable,colback=teal!4,colframe=teal!35!black,title={Before you move on}]
What does \kn{ಭ್ರಮೆ} mean? (\textquotedblleft{}An illusion\textquotedblright{}.) Say it once more.
\end{tcolorbox}
\section[\texorpdfstring{\kn{ಅಸ್ತಿತ್ವ} (astitva)}{ಅಸ್ತಿತ್ವ (astitva)}]{\kn{ಅಸ್ತಿತ್ವ} (astitva) — existence}

\label{lesson:KA-C382-astitva}



\begin{quote}
Before the new one: say the Kannada for reality, then the Kannada for an illusion.
\end{quote}

\subsection*{You'll want to know: \kn{ಅಸ್ತಿತ್ವ}}
\textbf{\kn{ಅಸ್ತಿತ್ವ}} — \emph{astitva} — \textquotedblleft{}existence\textquotedblright{}.

\begin{quote}
\kn{ಆ} \kn{ಹಳೆಯ} \kn{ಕೆರೆ} \kn{ಈಗ} \kn{ಅಸ್ತಿತ್ವದಲ್ಲಿ} \kn{ಇಲ್ಲ}.
\end{quote}

\begin{quote}
\emph{ā haḷeya kere īga astitvadalli illa.}
\end{quote}

\begin{quote}
That old lake no longer exists.
\end{quote}

\kn{ಅಸ್ತಿತ್ವದಲ್ಲಿ} \kn{ಇರು}, to exist; \kn{ಅಸ್ತಿತ್ವಕ್ಕೆ} \kn{ಬರು}, to come into being.

\begin{cousinweb}
\textbf{\kn{ಅಸ್ತಿತ್ವ}} (\emph{astitva}) is borrowed from Sanskrit \emph{astitva}, \textquotedblleft{}existence\textquotedblright{} (CDIAL 978), from \emph{asti}, \textquotedblleft{}it is\textquotedblright{}.
\end{cousinweb}

\begin{grammarlens}[title={What you've built}]
One more word for giving reasons and opinions.
\end{grammarlens}

\subsection*{Your turn}
\begin{itemize}
\raggedright
  \item \emph{Say it:} \emph{astitva}
  \item \emph{Say it:} \emph{astitva}, once more
  \item \emph{Say it:} \emph{ā haḷeya kere īga astitvadalli illa.} — That old lake no longer exists.
  \item \emph{Recall:} say the Kannada for reality, then the Kannada for an illusion, then say \emph{astitva} again
\end{itemize}

\begin{tcolorbox}[breakable,colback=teal!4,colframe=teal!35!black,title={Before you move on}]
What does \kn{ಅಸ್ತಿತ್ವ} mean? (\textquotedblleft{}Existence\textquotedblright{}.) Say it once more.
\end{tcolorbox}
\section[\texorpdfstring{\kn{ಮನಗಾಣು} (managāṇu)}{ಮನಗಾಣು (managāṇu)}]{\kn{ಮನಗಾಣು} (managāṇu) — to realise}

\label{lesson:KA-C382-managanu}



\begin{quote}
Before the new one: say the Kannada for an illusion, then the Kannada for existence.
\end{quote}

\subsection*{You'll want to know: \kn{ಮನಗಾಣು}}
\textbf{\kn{ಮನಗಾಣು}} — \emph{managāṇu} — \textquotedblleft{}to realise\textquotedblright{}.

\begin{quote}
\kn{ನನ್ನ} \kn{ತಪ್ಪನ್ನು} \kn{ನಾನು} \kn{ತಡವಾಗಿ} \kn{ಮನಗಂಡೆ}.
\end{quote}

\begin{quote}
\emph{nanna tappannu nānu taḍavāgi managaṇḍe.}
\end{quote}

\begin{quote}
I realised my mistake late.
\end{quote}

Past (I …): \kn{ಮನಗಂಡೆ}, \emph{managaṇḍe}.

Its past is built on \kn{ಕಂಡ}: \kn{ಮನಗಂಡನು}, he realised.

\begin{cousinweb}
\textbf{\kn{ಮನಗಾಣು}} (\emph{managāṇu}) is built from two parts: \emph{mana}, \textquotedblleft{}mind\textquotedblright{} (Sanskrit \emph{manas}, CDIAL 9822), and native \emph{kāṇ}, \textquotedblleft{}to see\textquotedblright{} (DEDR 1443): to see with the mind.
\end{cousinweb}

\begin{grammarlens}[title={What you've built}]
One more word for giving reasons and opinions.
\end{grammarlens}

\subsection*{Your turn}
\begin{itemize}
\raggedright
  \item \emph{Say it:} \emph{managāṇu}
  \item \emph{Say it:} \emph{managāṇu}, once more
  \item \emph{Say it:} \emph{nanna tappannu nānu taḍavāgi managaṇḍe.} — I realised my mistake late.
  \item \emph{Recall:} say the Kannada for an illusion, then the Kannada for existence, then say \emph{managāṇu} again
\end{itemize}

\begin{tcolorbox}[breakable,colback=teal!4,colframe=teal!35!black,title={Before you move on}]
What does \kn{ಮನಗಾಣು} mean? (\textquotedblleft{}To realise\textquotedblright{}.) Say it once more.
\end{tcolorbox}
